\section{Policy Model}

ActPlane policies are object-level information-flow rules over three node
types: processes, files, and endpoints. A process node represents a Linux task
identified by pid and command identity. A file node is currently keyed by path.
An endpoint node is currently keyed by IPv4 address. The model is intentionally
coarser than byte-level dynamic taint tracking~\cite{kemerlis2012libdft}; it
trades precision for low-overhead whole-system coverage.

\subsection{Labels and Propagation}

Each node carries a bitset of labels. A source declaration introduces labels:
an exec source labels processes that execute a matching binary, a file source
labels matching files, and an endpoint source labels matching network
endpoints. Labels propagate through fixed transfer functions:

\begin{itemize}
\item fork: a child process inherits the parent's labels;
\item exec: a process absorbs labels from the executed file and matching exec
sources;
\item read: a process absorbs labels from the file it reads;
\item write: a file absorbs labels from the process writing it;
\item connect: an endpoint records labels from the connecting process.
\end{itemize}

This is the provenance invariant ActPlane maintains: if a process carries label
\texttt{SECRET}, then information from some \texttt{SECRET} source has reached
that process through observed OS-level flow edges.

\subsection{Rules, Gates, and Effects}

A rule contains one or more deny clauses. A clause specifies an operation
(\texttt{exec}, \texttt{read}, \texttt{write}, \texttt{unlink},
\texttt{connect}, or \texttt{open}), a target pattern, a Boolean expression over
labels, and an optional condition such as a target allow-list, a lineage gate,
or an after gate.

\begin{figure}[t]
\begin{verbatim}
label AGENT
source SECRET = file "**/.env"

rule no-secret-egress:
  deny connect endpoint "*" if SECRET
  effect block
  reason "secret-derived data must stay local"
  remediation "run the redactor or ask the user"

rule no-agent-git:
  deny exec "**/git" if AGENT
  effect kill
  reason "agent runs must not invoke git"
\end{verbatim}
\caption{A shortened ActPlane policy. The first rule asks for LSM-only
pre-operation deny. The second rule asks for harness-level termination and can
fire from tracepoint observation.}
\Description{A code listing showing two ActPlane rules: one blocking secret
egress and one killing git execution by an agent.}
\label{fig:policy-example}
\end{figure}

Rules have explicit effects.

\textbf{Audit.} The kernel reports the violation and allows the action to
proceed. This is useful for soft workflow guidance and for measurement.

\textbf{Block.} A BPF-LSM hook returns \texttt{-EPERM} before the operation is
committed. \texttt{block} is not supported in tracepoint mode and is not
silently downgraded to \texttt{audit} or \texttt{kill}.

\textbf{Kill.} The kernel sends \texttt{SIGKILL} to the current task. If a
pre-operation LSM hook is active, ActPlane can also return \texttt{-EPERM};
from a tracepoint, \texttt{kill} is harness-level termination rather than
pre-operation denial.

If several supported rules match the same event, ActPlane chooses the strongest
effect in the order \texttt{kill > block > audit}. Unsupported effects are
ignored by that backend, so a tracepoint \texttt{audit} rule can still fire even
when an unsupported \texttt{block} rule also matches.

\subsection{Why This Is a Harness DSL}

The language is deliberately not a general-purpose MAC policy. Its primitives
map to agent workflows: ``after tests,'' ``through a review tool,'' ``data
derived from untrusted input,'' and ``read-only sub-agent.'' The grammar is
small because policy authors should be able to translate natural-language
instructions into a few recurring source/sink/gate shapes.
